\documentclass[aspectratio=169,11pt]{beamer}

% ---------------------------------------------------------------
% Packages
% ---------------------------------------------------------------
\usepackage[T1]{fontenc}
\usepackage{lmodern}
\usepackage{fontawesome5}
\usepackage{tabularx,booktabs}
\usepackage{listings}
\usepackage{dirtree}
\usepackage[most]{tcolorbox}

% ---------------------------------------------------------------
% Design: minimal light blue
% ---------------------------------------------------------------
\definecolor{accent}{RGB}{45,110,170}
\definecolor{rulecol}{RGB}{160,195,230}
\definecolor{boxfill}{RGB}{232,241,250}
\definecolor{muted}{RGB}{110,110,110}

\setbeamertemplate{navigation symbols}{}
\setbeamercolor{frametitle}{fg=accent}
\setbeamercolor{title}{fg=accent}
\setbeamercolor{subtitle}{fg=muted}
\setbeamercolor{structure}{fg=accent}
\setbeamercolor{normal text}{fg=black!85}
\setbeamerfont{frametitle}{size=\Large}
\setbeamertemplate{itemize items}{\textcolor{accent}{\tiny$\blacksquare$}}
\setbeamertemplate{enumerate items}[default]

\setbeamertemplate{frametitle}{%
  \vspace*{0.5em}%
  \usebeamerfont{frametitle}\usebeamercolor[fg]{frametitle}\insertframetitle\par
  \vspace{-0.3em}{\color{rulecol}\rule{\linewidth}{0.6pt}}%
}
\setbeamertemplate{footline}{%
  \hfill{\color{muted}\footnotesize\insertframenumber}\hspace*{1.2em}\vspace*{0.8em}%
}

% Prompt boxes: example instructions to the agent
\newtcolorbox{promptbox}{colback=boxfill, colframe=boxfill, boxrule=0pt,
  arc=2pt, left=5pt, right=5pt, top=2pt, bottom=2pt, before skip=3pt,
  after skip=3pt, fontupper=\small\itshape\raggedright}
\newcommand{\prompt}[1]{\begin{promptbox}\textcolor{accent}{\faCommentDots}\enspace #1\end{promptbox}}

% Highlight boxes: key messages
\newtcolorbox{keybox}{colback=white, colframe=rulecol, boxrule=0.6pt,
  arc=2pt, left=6pt, right=6pt, top=3pt, bottom=3pt}

% Code listings (examples are included from files)
\lstset{basicstyle=\ttfamily\scriptsize, backgroundcolor=\color{boxfill},
  frame=none, columns=fullflexible, keepspaces=true, upquote=true,
  breaklines=true, aboveskip=3pt, belowskip=3pt, xleftmargin=3pt,
  xrightmargin=3pt, commentstyle=\color{muted}}

% Directory trees
\renewcommand{\DTstyle}{\small}
\renewcommand{\DTstylecomment}{\color{muted}\footnotesize}
\setlength{\DTbaselineskip}{12pt}
\DTsetlength{0.2em}{1em}{0.2em}{0.4pt}{1.6pt}

\newcolumntype{Y}{>{\raggedright\arraybackslash}X}
\newcommand{\file}[1]{\texttt{#1}}
\newcommand{\email}{\mbox{\href{mailto:simon.martin@univie.ac.at}{simon.martin@univie.ac.at}}}
\newcommand{\icon}[1]{\textcolor{accent}{#1}}
\newcommand{\repourl}{\href{https://github.com/simonmartinecon/jumpstart-ai}{github.com/simonmartinecon/jumpstart-ai}}

% ---------------------------------------------------------------
% Title
% ---------------------------------------------------------------
\title{Jump-starting AI agents}
\subtitle{A practical quick start for economists}
\author{Simon Martin}
\institute{\footnotesize
  \faEnvelope\enspace \email\\[2pt]
  \faGlobe\enspace \href{https://sites.google.com/view/simonmartinecon/}{sites.google.com/view/simonmartinecon}\\[2pt]
  \faGithub\enspace \repourl}
\date{\footnotesize Version 0.1, October 2026}
\newcommand{\disclaimer}{These slides focus on practical aspects. There is, of course, much discussion about the ethics of AI use (in writing and reviewing), data policies and more; assessing these is up to you.}

\begin{document}

% ---------------------------------------------------------------
\begin{frame}[plain]
  \titlepage
  \vfill
  {\color{rulecol}\rule{\linewidth}{0.4pt}}\par
  {\scriptsize\color{muted}\disclaimer\par}
\end{frame}

% ---------------------------------------------------------------
\section{Introduction}

\begin{frame}{About these slides}
  \begin{itemize}
    \item \textbf{Objective:} get you started on one of your own projects right after skimming these slides.
    \item \textbf{For whom:} economists in academia, but much of this applies more widely.
    \item \textbf{Prerequisites:} almost none. To start on your project right away, you need a basic subscription (about 20~USD/month) with any AI provider.
    \item \textbf{Caveat:} this is my personal take in a field that changes very fast. It is what works for me, not a fixed recipe: adjust it to your own workflow and software.
    \item \textbf{Meta:} these slides were made with an agent, following exactly the workflow they describe. The full project folder is on GitHub (\repourl).
  \end{itemize}
  \vspace{0.5em}
  \begin{keybox}
    \icon{\faComments}\enspace Happy to discuss with anyone, any time:\enspace \email
  \end{keybox}
\end{frame}

% ---------------------------------------------------------------
\begin{frame}{Chatbot vs agent}
  \renewcommand{\arraystretch}{1.35}
  \begin{tabularx}{\linewidth}{@{}l Y Y@{}}
    & \icon{\faCommentDots}\enspace\textbf{Chatbot} & \icon{\faUserTie}\enspace\textbf{Agent} \\
    \midrule
    Input   & you copy/paste text into the chat & reads your files directly \\
    Output  & you copy/paste answers back into your documents & writes into your files (text, code, tables) \\
    Context & the current chat only & your whole project folder \\
    Tasks   & answers one question at a time & plans and carries out multi-step tasks on its own \\
    \bottomrule
  \end{tabularx}
  \vspace{1em}
  \begin{keybox}
    The basics on these slides get you maybe 90--95\% of the gains. Ignore advanced features (skills, GitHub integration, \dots) for now: start simple.
  \end{keybox}
\end{frame}

% ---------------------------------------------------------------
\begin{frame}{``Coding agent'': not just for coding}
  \begin{itemize}
    \item \textbf{Misconception 1: coding is for programmers.} Almost everything you do on a computer is ``coding'': writing and formatting text, reading documents, calculating, combining content from several files.
    \item \textbf{Misconception 2: it is for empirical work only.} Agents work just as well for writing, admin and theory; their maths skills have improved a lot recently.
  \end{itemize}
  \vspace{0.3em}
  \begin{columns}[T]
    \begin{column}{0.48\textwidth}
      \icon{\faPen}\enspace\textbf{Writing}
      \prompt{Check that the introduction of paper.tex is consistent with Section 5.}
      \icon{\faClipboardList}\enspace\textbf{Admin}
      \prompt{Draft the progress report from the documents in funding/.}
    \end{column}
    \begin{column}{0.48\textwidth}
      \icon{\faInfinity}\enspace\textbf{Theory}
      \prompt{Turn my notes in notes.md into a formal proposition, in the style of the samples in samples/.}
      \icon{\faChartBar}\enspace\textbf{Empirics}
      \prompt{Add clustered standard errors to all regressions in code/analysis.do.}
    \end{column}
  \end{columns}
\end{frame}

% ---------------------------------------------------------------
\begin{frame}{What it is, and what it isn't}
  \begin{columns}[T]
    \begin{column}{0.48\textwidth}
      \icon{\faQuestionCircle}\enspace\textbf{``Just a token-guessing machine?''}
      \begin{itemize}
        \item Technically, yes: it predicts the next piece of text.
        \item But predicting well at this scale means it can follow arguments, plan tasks and check proofs.
        \item Judge it by its results on your own tasks.
      \end{itemize}
    \end{column}
    \begin{column}{0.48\textwidth}
      \icon{\faDatabase}\enspace\textbf{Not a store of truth}
      \begin{itemize}
        \item It is not a database, and not Wikipedia.
        \item It reconstructs facts rather than looking them up, so they can be wrong, outdated or invented (citations!).
      \end{itemize}
    \end{column}
  \end{columns}
  \vspace{1em}
  \begin{keybox}
    \icon{\faLightbulb}\enspace Strong at working with what you give it, weak at recalling facts. So provide the sources (literature/, data/), and let it compute in code instead of answering from memory (see slides~\ref{script} and~\ref{pitfalls}).
  \end{keybox}
\end{frame}

% ---------------------------------------------------------------
\begin{frame}{What you need}
  \begin{columns}[T]
    \begin{column}{0.48\textwidth}
      \icon{\faTools}\enspace\textbf{Getting started}
      \begin{itemize}\small
        \item \textbf{Subscription:} any basic plan at about 20~USD/month (Claude, ChatGPT/Codex, Gemini, \dots). Free plans are usually too limited; 200~USD plans are rarely needed. The provider matters less than you might think.
        \item \textbf{Software:} an editor with the agent built in, e.g.\ VS Code with the Claude Code extension (my setup).
        \item \textbf{Setup:} ask the agent to install or explain LaTeX, Stata paths etc., and to convert existing .docx/.pdf drafts to .tex/.md once.
      \end{itemize}
    \end{column}
    \begin{column}{0.48\textwidth}
      \icon{\faSlidersH}\enspace\textbf{Settings}
      \begin{itemize}\small
        \item \textbf{Model:} any current model is fine. No need to chase every new release.
        \item \textbf{Effort:} medium for most tasks; it trades depth against speed and tokens. High only for occasional thorough checks.
        \item \textbf{Tokens:} the 20~USD plan suffices for most daily use. What eats tokens: many parallel web searches, and binary files (.docx, .xlsx, .pdf) instead of plain text.
      \end{itemize}
    \end{column}
  \end{columns}
\end{frame}

% ---------------------------------------------------------------
\section{Core mindset}

\begin{frame}{Your agent is your research assistant}
  \begin{itemize}
    \item Extremely capable, but only as good as the context you give it.
    \item Your job, as with any new RA:
    \begin{itemize}
      \item \textbf{Context:} background, related literature, your terminology. A brilliant RA with a physics background needs to know how you use certain terms.
      \item \textbf{Instructions:} what to do, which output you expect, and in which file.
      \item \textbf{Unclear points:} say how to handle them, for example:
    \end{itemize}
  \end{itemize}
  \prompt{Never make assumptions. Always clarify with me first.}
  \prompt{Proceed as far as you can. Leave a TODO note in the document for anything unclear, so we can come back to it later.}
  \begin{itemize}
    \item You do most of this once, in \file{instructions.md} (see slide~\ref{instructions}).
  \end{itemize}
\end{frame}

% ---------------------------------------------------------------
\begin{frame}{You stay in charge}
  \begin{itemize}
    \item RAs make mistakes and misread input, and so does your agent. \textbf{You remain responsible for the final output.}
    \item \textbf{Verify:} read and check everything, and leave enough time for it at the end. AI slop shows.
    \item A second agent can cross-check for you, but the responsibility stays with you.
    \item \textbf{Memory drifts:} as a session grows, older or partly contradictory instructions get lost. Don't assume everything you ever said is still respected; now and then ask it to re-read the instructions or double-check.
    \item \textbf{Fresh sessions:} start new sessions rather than one endless chat. Let the agent record progress and open points in \file{status.md}, and read it at the start of each session.
  \end{itemize}
\end{frame}

% ---------------------------------------------------------------
\begin{frame}{Project = folder}
  \begin{itemize}
    \item Agents work in \textbf{projects}, and a project is simply a folder on your computer.
    \item Everything the agent needs goes in there, and you come back to the same folder each time. Many projects are short-lived; that is fine.
  \end{itemize}
  \vspace{0.3em}
  \begin{columns}[T]
    \begin{column}{0.48\textwidth}
      \textbf{A research paper}\par\vspace{0.3em}
      \dirtree{%
.1 \faFolderOpen\ my-paper/.
.2 instructions.md\DTcomment{the contract}.
.2 status.md\DTcomment{progress log}.
.2 paper.tex\DTcomment{main document}.
.2 \faFolder\ literature/\DTcomment{PDFs you may cite}.
.2 \faFolder\ samples/\DTcomment{your writing samples}.
.2 \faFolder\ data/\DTcomment{raw data}.
.2 \faFolder\ code/\DTcomment{.do files}.
.2 \faFolder\ output/\DTcomment{generated tables, figures}.
}

    \end{column}
    \begin{column}{0.48\textwidth}
      \textbf{These slides}\par\vspace{0.3em}
      \dirtree{%
.1 \faFolderOpen\ jumpstart-slides/.
.2 instructions.md\DTcomment{the contract}.
.2 status.md\DTcomment{progress log}.
.2 background.md\DTcomment{my notes}.
.2 \faFolder\ recordings/\DTcomment{voice transcripts}.
.2 structure.md\DTcomment{agreed outline}.
.2 \faFolder\ examples/\DTcomment{code shown on slides}.
.2 slides.tex\DTcomment{main document}.
.2 slides.pdf\DTcomment{final output}.
}

    \end{column}
  \end{columns}
\end{frame}

% ---------------------------------------------------------------
\begin{frame}{Plain text for everything you work on}
  \renewcommand{\arraystretch}{1.3}
  \begin{tabularx}{\linewidth}{@{}Y Y@{}}
    \icon{\faFile*}\enspace\textbf{Plain text} & \icon{\faFileArchive}\enspace\textbf{Binary} \\
    \midrule
    .tex, .md, .txt, .csv, .do, .R, .py & .docx, .xlsx, .pdf, .dta, .png \\
    $\rightarrow$ everything you actively work on & $\rightarrow$ only pure input (literature, external files) and final output (the compiled PDF) \\
    \bottomrule
  \end{tabularx}
  \vspace{0.3em}
  \begin{itemize}
    \item \textbf{Why:} plain text is close to how the agent reads: much faster, far fewer tokens, and much easier to edit.
    \item \textbf{Example:} to find GDP for Austria in 2020, the agent first has to decode an .xlsx file. In a .csv, one quick search returns the number.
    \item \textbf{Bonus:} changes over time are easy to see (what changed, when, and why), even with simple Dropbox file history.
  \end{itemize}
  \begin{keybox}
    \icon{\faLightbulb}\enspace\textbf{Try it:} open a .docx or .pdf in Notepad or TextEdit. That is what the agent has to work through.
  \end{keybox}
\end{frame}

% ---------------------------------------------------------------
\begin{frame}{No copy/paste}
  \begin{itemize}
    \item Copy/paste is friction: you should never need it.
    \item \textbf{Input:} point the agent to where things are.
  \end{itemize}
  \prompt{Consider Smith2023.pdf in literature/.}
  \prompt{Rework paragraph 3 of the introduction in paper.tex.}
  \prompt{Check the equation in line 142 of paper.tex.}
  \begin{itemize}
    \item \textbf{Output:} have it write straight into the final file. It finds the right place.
  \end{itemize}
  \prompt{Slide 17 is fine as discussed. Write it into slides.tex.}
  \begin{itemize}
    \item Never copy answers from the chat window into your documents.
  \end{itemize}
\end{frame}

% ---------------------------------------------------------------
\begin{frame}[fragile]{Script everything}
  \label{script}
  \begin{itemize}\small
    \item Every table, figure and number in your document comes from a script. Never edit generated tables by hand: you will redo each one ten more times, so make it come out right from the code.
    \item You don't need to know how to script an operation: the agent does it for you.
  \end{itemize}
  \begin{columns}[T]
    \begin{column}{0.52\textwidth}
      \lstset{basicstyle=\ttfamily\tiny}
      \footnotesize\textbf{code/wage\_regs.do} writes the table:
      \lstinputlisting{examples/wage_regs.do}
      \textbf{paper.tex} includes it:
      \lstinputlisting{examples/paper_input.tex}
    \end{column}
    \begin{column}{0.44\textwidth}
      \footnotesize
      \renewcommand{\arraystretch}{1.3}
      \begin{tabularx}{\linewidth}{@{}Y Y@{}}
        \toprule
        \textcolor{red!70!black}{\faTimes}\enspace\textbf{Don't} & \textcolor{green!50!black}{\faCheck}\enspace\textbf{Do} \\
        \midrule
        Change the column order of Table~2 in paper.tex. & Change wage\_regs.do so that Table~2 comes out in this column order. \\
        \midrule
        What does this coefficient imply in percent? (Then paste the answer.) & Add this calculation to the do-file, export the number, and use it in Section~4. \\
        \bottomrule
      \end{tabularx}
    \end{column}
  \end{columns}
\end{frame}

% ---------------------------------------------------------------
\section{instructions.md}

\begin{frame}{\file{instructions.md}: the contract}
  \label{instructions}
  \begin{columns}[T]
    \begin{column}{0.6\textwidth}
      \footnotesize
      \renewcommand{\arraystretch}{1.2}
      \begin{tabularx}{\linewidth}{@{}l Y@{}}
        \toprule
        \textbf{Section} & \textbf{My typical content} \\
        \midrule
        Scope      & \textit{We are writing a research paper in labour economics.} \\
        Literature & \textit{Read the papers in literature/.} \\
        Stages     & sub-steps, and how we proceed and interact (longer projects) \\
        Output     & \textit{The main document is paper.tex.} \\
        Software   & which software; where code, input and output files are \\
        Style      & a writing sample of mine in samples/: sets layout, language and tone \\
        Checks     & \textit{At the end: grammar, language, style and consistency check.} \\
        Logging    & \textit{Record progress and TODOs in status.md; read it at the start of each session.} \\
        \bottomrule
      \end{tabularx}
    \end{column}
    \begin{column}{0.36\textwidth}
      \begin{itemize}\small
        \item Instead of long one-off prompts: one file that sets the scope and the rules.
        \item \textbf{Format barely matters:} no syntax rules, .md or .txt. A little structure (headings) helps.
        \item \textbf{Why bother:} 10--15 minutes of writing forces you to think through background, deliverables and sub-steps, and tells your RA clearly what to do.
        \item Extend it as you go.
      \end{itemize}
    \end{column}
  \end{columns}
\end{frame}

% ---------------------------------------------------------------
\begin{frame}[fragile]{Example: the \file{instructions.md} behind these slides}
  \lstset{basicstyle=\ttfamily\fontsize{5.5}{6.6}\selectfont}
  \begin{columns}[T]
    \begin{column}{0.49\textwidth}
      \lstinputlisting[linerange={1-10}]{instructions.md}
    \end{column}
    \begin{column}{0.49\textwidth}
      \lstinputlisting[linerange={12-20}]{instructions.md}
    \end{column}
  \end{columns}
  \vspace{0.5em}
  {\fontsize{8}{10}\selectfont\icon{\faGithub}\enspace Excerpt. The full file is on GitHub: \href{https://github.com/simonmartinecon/jumpstart-ai/blob/main/instructions.md}{github.com/simonmartinecon/jumpstart-ai/blob/main/instructions.md}\par}
\end{frame}

% ---------------------------------------------------------------
\section{Ways of working}

\begin{frame}{Ways of working}
  \begin{columns}[T]
    \begin{column}{0.48\textwidth}
      \icon{\faClock}\enspace\textbf{Asynchronous}
      \begin{itemize}\small
        \item Longer tasks run in the background: half an hour, a few hours, overnight. Run several agents on different projects; reconcile when they finish.
        \item Settle all decisions up front, or defer non-critical ones as TODOs. What to avoid: after an hour, the agent stops to ask something you could have settled at the start.
      \end{itemize}
      \prompt{Plan the subtasks for this stage and ask me about anything unclear before you start. Leave non-critical points as TODOs.}
    \end{column}
    \begin{column}{0.48\textwidth}
      \icon{\faComments}\enspace\textbf{Sparring partner}
      \begin{itemize}\small
        \item Test ideas before taking them to colleagues: does the argument hold up against what we wrote elsewhere, and against the literature?
        \item Often points out what you have overlooked, such as a contradiction with an earlier section.
      \end{itemize}
      \prompt{I think we should reframe the discussion in Section 5 as follows: \dots{} Assess this critically against the rest of the paper and the literature in literature/.}
    \end{column}
  \end{columns}
  \vspace{0.2em}
  {\small Working interactively on one part is fine too, but most tasks can run asynchronously.}
\end{frame}

% ---------------------------------------------------------------
\begin{frame}{Pitfalls}
  \label{pitfalls}
  \begin{itemize}
    \item \icon{\faBook}\enspace\textbf{Made-up or outdated citations.} My rule in \file{instructions.md}:
  \end{itemize}
  \prompt{NEVER cite (add to the .bib file) any document that is not in literature/.}
  \begin{itemize}
    \item[] The agent may suggest literature, but I look up and download the papers myself. This way I have the real, latest versions, and cite only what I chose.
    \item \icon{\faSave}\enspace\textbf{Backups.} The agent writes into your files. Keep file history (Dropbox, OneDrive, \dots), and read permission prompts before approving them.
    \item \icon{\faLock}\enspace\textbf{Confidential data.} Check data licences and your institution's rules before giving the agent access to confidential data.
  \end{itemize}
\end{frame}

% ---------------------------------------------------------------
\section{Jump start}

\begin{frame}{Putting it together: the jump start}
  \begin{columns}[T]
    \begin{column}{0.5\textwidth}
      \begin{enumerate}
        \item Create a folder for the project.
        \item Put in the context files: literature, data, a writing sample, your draft or notes.
        \item Write \file{instructions.md} (10--15 minutes).
        \item Open the agent in that folder (e.g.\ VS Code with Claude Code).
        \item Tell it to read the instructions, and take it from there.
      \end{enumerate}
    \end{column}
    \begin{column}{0.46\textwidth}
      \textbf{First prompts}
      \prompt{Read instructions.md. Plan the subtasks and ask me about anything unclear before you start.}
      \prompt{Read status.md and continue where we left off.}
      \prompt{Do a thorough grammar, style and consistency check of paper.tex.}
    \end{column}
  \end{columns}
\end{frame}

% ---------------------------------------------------------------
\begin{frame}{Get started now}
  \begin{columns}[T]
    \begin{column}{0.22\textwidth}
      \centering
      \parbox[c][1.3cm][c]{\linewidth}{\centering\Huge\icon{\faFolderPlus}}\par
      \textbf{1. Create a folder}\par\smallskip
      {\small one per project}
    \end{column}
    \begin{column}{0.22\textwidth}
      \centering
      \parbox[c][1.3cm][c]{\linewidth}{\centering\Huge\icon{\faCopy}}\par
      \textbf{2. Put in the context files}\par\smallskip
      {\small literature, data, notes, a writing sample}
    \end{column}
    \begin{column}{0.22\textwidth}
      \centering
      \parbox[c][1.3cm][c]{\linewidth}{\centering\Huge\icon{\faFileSignature}}\par
      \textbf{3. Write \file{instructions.md}}\par\smallskip
      {\small scope, output, stages, style}
    \end{column}
    \begin{column}{0.22\textwidth}
      \centering
      \parbox[c][1.3cm][c]{\linewidth}{\centering\Huge\icon{\faPlay}}\par
      \textbf{4. Open the agent and prompt:}
    \end{column}
  \end{columns}
  \vspace{1.5em}
  \begin{center}
    \begin{minipage}{0.5\textwidth}
      \prompt{\normalsize Read the instructions.}
    \end{minipage}
  \end{center}
\end{frame}

% ---------------------------------------------------------------
\begin{frame}{Thanks and further reading}
  \begin{itemize}
    \item \textbf{Anton Korinek:} big thanks! A great source on using AI; I have benefited a lot from his talks and from interacting with him. His series of JEL articles on AI for economic research (latest: ``AI Agents for Economic Research'', 2025) is at \href{https://www.korinek.com/Home}{korinek.com}.
    \item \textbf{Scott Cunningham:} Scott's Mixtape Substack, on causal inference and how AI agents change empirical research: \href{https://causalinf.substack.com}{causalinf.substack.com}.
    \item \textbf{Paul Goldsmith-Pinkham:} great blog, Substack and video series: \href{https://paulgp.com/blog.html}{paulgp.com/blog.html}.
  \end{itemize}
  \vspace{0.8em}
  \begin{keybox}
    \icon{\faComments}\enspace Questions, or a workflow that works better for you? Happy to discuss: \email
  \end{keybox}
\end{frame}

\end{document}
